\subsection{Methods, Algorithms, and Computational Artifacts}

Procedural blocks separate conceptual guidance, algorithm specifications, implementations, and correctness or performance evidence.

\begin{convention} \label{conv:procedural-environment-roles} \hlt{Procedural Environment Roles}\\
Use the following procedural roles throughout these notes.
\begin{enumerate}[label=(\roman*)]
\setlength{\itemsep}{0pt}
\item \hlt{Proof method}: a reusable proof route, reduction, recognition criterion, or obligation checklist.
\item \hlt{Computation method}: a reusable construction, calculation, approximation, diagnostic, or test.
\item \hlt{Method}: a short convention, classification workflow, or inseparable hybrid with no clearer subtype.
\item \hlt{Process}: an organizational workflow for research, writing, review, or documentation.
\item \hlt{Algorithm}: executable pseudocode with declared inputs, outputs, control flow, and termination behavior when applicable.
\end{enumerate}
Worked instances belong in examples, completed arguments in proofs, and implementation-specific behavior in code or artifact records.
\end{convention}

\begin{process} \label{proc:algorithm-specification-audit} \hlt{Algorithm Specification and Audit}\\
Before publishing pseudocode, record its contract and associated verification obligations.
\begin{enumerate}[label=(\roman*)]
\setlength{\itemsep}{0pt}
\item \hlt{Problem}: state the task and the abstraction level at which the algorithm operates.
\item \hlt{Input}: specify types, dimensions, admissibility, representation, and malformed-input behavior.
\item \hlt{Output}: specify the returned object and a postcondition that characterizes success.
\item \hlt{State}: define initialization, mutable state, determinism or randomness, and any invariant used for correctness.
\item \hlt{Steps}: write one unambiguous action per step with readable branches and loops.
\item \hlt{Termination}: state the stopping rule and prove termination when it is claimed.
\item \hlt{Correctness}: state the proof obligations and keep the proof of the postcondition separate from the pseudocode.
\item \hlt{Resources}: state complexity only with an input-size parameter, cost model, resource, and case type.
\item \hlt{Limits}: record numerical, memory, precision, platform, approximation, and applicability boundaries.
\end{enumerate}
Labeled steps and explanatory invariants clarify the distinction between an algorithm and program text \citep{knuth1988mathematicalwriting}.
\end{process}

\begin{principle} \label{princ:formal-informal-procedural-pair} \hlt{Formal and Informal Procedural Views}\\
Use complementary formal and informal views when a procedure is complex. The formal view fixes inputs, outputs, updates, and stopping conditions; the informal view explains the invariant, decomposition, or design choice. The views should not duplicate one another, and only a separate proof can establish mathematical correctness.
\end{principle}

\begin{process} \label{proc:computational-artifact-record} \hlt{Computational Artifact Record}\\
For code supporting a reported result, preserve an auditable artifact record.
\begin{enumerate}[label=(\roman*)]
\setlength{\itemsep}{0pt}
\item \hlt{Identity}: archive the exact source revision and data identifiers.
\item \hlt{Environment}: record dependencies, runtime, platform, configuration, and random seeds.
\item \hlt{Access}: state licenses, unavailable components, and access restrictions.
\item \hlt{Execution}: record the commands and expected outputs.
\item \hlt{Traceability}: map each reported figure, table, or metric to its generating command.
\end{enumerate}
The 2026 NeurIPS checklist requests exact execution instructions and sufficient experimental detail for interpreting results \citep{neurips2026checklist}.
\end{process}

\begin{warning} \label{warn:algorithm-evidence-boundaries} \hlt{Algorithm Evidence Boundaries}\\
Passing tests does not establish correctness outside the tested cases, and an asymptotic bound does not predict wall-clock performance without a cost model and operating conditions. Keep specification, proof, implementation, benchmark, and interpretation distinct, and restrict each claim to the evidence that supports it.
\end{warning}
